\documentclass[runningheads]{llncs}

\usepackage[T1]{fontenc}
\usepackage{graphicx}
\usepackage{amsmath,amssymb,amsfonts}
\usepackage{booktabs}
\usepackage{multirow}
\usepackage{xcolor}

\begin{document}

\title{VD-TTFS: Efficient TTFS Computing Architecture on Vector Devices\\
       \large Supplementary Material}
\titlerunning{VD-TTFS (Supplementary Material)}

\author{Shifeng Liu\inst{1} \and
Zhijie Yang\inst{1} \and
Dongsheng Wang\inst{1} \and
Lei Wang\inst{1}}

\authorrunning{S. Liu et al.}

\institute{Defense Innovation Institute, Academy of Military Sciences, Beijing, China\\
\email{charlley.g.l@gmail.com}}

\maketitle

% ---- switch to appendix numbering (A, B, C, ...) ----
\appendix

\section{Selecting the Integrator Chunk Width $\Delta$}\label{app:chunk}

The time-sorted early-exit integrator of the main paper exposes a single
granularity hyper-parameter, the \emph{chunk width} $\Delta$: the number of
consecutive timesteps grouped into each block that the integrator sweeps. This
appendix details how $\Delta$ is chosen and reports the measured basis for the
per-task values used in the paper ($\Delta=8$ for VGG-16, $\Delta=16$ for LeNet,
and $\Delta=32$ for the DVSGesture network). All measurements were taken on a
single NVIDIA GeForce RTX~5070~Ti (16\,GB, CUDA~12.8, \texttt{nvcc -O3}); the
optima are hardware-specific, but the underlying relationship is not (Section~\ref{app:chunk-find}).

\subsection{Method}\label{app:chunk-method}

\paragraph{The trade-off.}
The integrator partitions the (truncated) window $[0,T'_\ell)$ into chunks of
$\Delta$ timesteps, accumulates the synaptic events arriving within each chunk
into a per-neuron register array \texttt{delta}$[\Delta]$, integrates chunk by
chunk, and \emph{halts the instant a neuron crosses its threshold}, so the events
of all later chunks are never fetched. $\Delta$ sets the resolution of this early
exit and induces two opposing effects:
\begin{itemize}
  \item A \textbf{finer} $\Delta$ lets a neuron stop at a finer temporal
    granularity, so fewer \emph{post-firing} synapses are processed---valuable
    when the arrival-ordered early exit is active. However, each chunk carries a
    fixed bookkeeping cost (clearing \texttt{delta}$[\Delta]$, a threshold scan,
    loop overhead), repeated $\lceil T/\Delta\rceil$ times, so a very small
    $\Delta$ multiplies that overhead.
  \item A \textbf{coarser} $\Delta$ issues fewer chunks and less bookkeeping, but
    blunts the early exit (a fired neuron still pays out its current chunk) and
    enlarges \texttt{delta}$[\Delta]$; beyond the register budget this array
    spills into local memory, which is slow.
\end{itemize}
The optimum therefore depends on how much early exit there is to exploit---i.e.\
the arrival-ordered early-exit factor $\gamma$---and on the window length $T$:
a strong early exit (small $\gamma$) rewards a fine $\Delta$, whereas a weak or
absent early exit ($\gamma\!\to\!1$) turns fine chunking into pure overhead and
favours a coarser $\Delta$, up to the point where \texttt{delta}$[\Delta]$
local-memory pressure dominates.

\paragraph{Determination procedure.}
For each task we sweep $\Delta$ over powers of two up to the window length,
recompile, run the \emph{full} evaluation set on the target GPU, and select the
latency minimum. Each configuration is executed three times and the minimum
inference time is reported to suppress scheduling jitter. The task accuracy is
recorded at every $\Delta$: re-chunking is an exact re-ordering of the same
synaptic events, so the prediction must be invariant, which we verify empirically
(the lone exception, due to floating-point summation order, is discussed in
Section~\ref{app:chunk-find}).

\subsection{Measured Sweeps}\label{app:chunk-data}

Tables~\ref{tab:chunk-t1}--\ref{tab:chunk-t3} report the sweep for each task. The
row marked $\star$ is the latency minimum; the shipped default is shown in bold.

\begin{table}[t]
\centering
\caption{Task~1, LeNet/MNIST ($T=80$), full 10{,}000-image test set.}
\label{tab:chunk-t1}
\begin{tabular}{cccc}
\toprule
$\Delta$ & Accuracy & Inference time (s) & \\
\midrule
4  & 98.18\% & 0.1025 & \\
8  & 98.18\% & 0.0645 & \\
\textbf{16} & \textbf{98.18\%} & \textbf{0.0450} & $\star$ \\
32 & 98.18\% & 0.0489 & \\
64 & 98.18\% & 0.0559 & \\
\bottomrule
\end{tabular}
\end{table}

\begin{table}[t]
\centering
\caption{Task~2, VGG-16/CIFAR-10 ($T=80$), full 10{,}000-image test set.}
\label{tab:chunk-t2}
\begin{tabular}{cccc}
\toprule
$\Delta$ & Accuracy & Inference time (s) & \\
\midrule
4  & 90.76\% & 2.730 & $\star$ \\
\textbf{8}  & \textbf{90.76\%} & \textbf{2.747} & ($+0.6\%$) \\
16 & 90.76\% & 3.120 & \\
32 & 90.76\% & 3.204 & \\
64 & 90.76\% & 3.427 & \\
\bottomrule
\end{tabular}
\end{table}

\begin{table}[t]
\centering
\caption{Task~3, DVSGesture ($T=160$), 1078 samples.}
\label{tab:chunk-t3}
\begin{tabular}{cccc}
\toprule
$\Delta$ & Accuracy & Inference time (ms) & \\
\midrule
8   & 96.29\% & 1027 & \\
16  & 96.29\% & 955  & $\star$ \\
\textbf{32}  & \textbf{96.29\%} & \textbf{966}  & ($+1.2\%$) \\
64  & 96.38\%$^{\dagger}$ & 1014 & \\
160 & 96.29\% & 1466 & \\
\bottomrule
\end{tabular}
\end{table}

\subsection{Findings}\label{app:chunk-find}

\begin{enumerate}
  \item \textbf{The shipped defaults are at or near the per-task optimum.}
    Task~1 ($\Delta=16$) is the exact minimum; Task~2 ($\Delta=8$) lies within
    $0.6\%$ of the minimum (the marginally faster $\Delta=4$ is within run-to-run
    noise); Task~3 ($\Delta=32$) is within $1.2\%$ of the minimum and is preferred
    for its warp alignment.
  \item \textbf{The optimum tracks $\gamma$ exactly as argued in the main paper.}
    The strong early exit of VGG-16 ($\gamma=0.62$) rewards fine chunks
    ($\Delta=4$--$8$); the weak early exit of LeNet ($\gamma\approx0.97$) makes a
    medium $\Delta=16$ best, as fine chunking is overhead and coarse chunking
    blunts the little early exit available; the early-exit-free, long-window
    DVSGesture stream ($\gamma\approx0.99$, $T=160$) favours a coarse
    $\Delta=16$--$32$, while the extreme $\Delta=160$ spills
    \texttt{delta}$[\Delta]$ to local memory and is the slowest configuration.
  \item \textbf{Chunking is lossless.} Accuracy is invariant across $\Delta$ for
    Tasks~1--2 and for Task~3 except at $\Delta=64$ ($^{\dagger}$), where it
    differs by a single sample ($96.38\%$ vs.\ $96.29\%$, i.e.\ $1039$ vs.\ $1038$
    of $1078$). This reflects IEEE-754 accumulation order, not an algorithmic
    change: different chunk widths sum the identical set of contributions in a
    different order, and one borderline neuron's first-spike time flips. The event
    set and its contributions are unchanged---a lossless permutation.
  \item \textbf{Hardware dependence.} The reported optima are specific to the
    RTX~5070~Ti; the precise minimum can shift by one step on other GPUs, but the
    qualitative rule---$\gamma$ and $T$ determine whether a fine or coarse $\Delta$
    is best---is hardware-independent.
\end{enumerate}

\section{A Fused Multi-Step LIF Lower Bound for SpikingJelly}\label{app:sjfused}

\subsection{Motivation}\label{app:sjfused-purpose}
SpikingJelly does not provide an efficient TTFS kernel, and the SpikingJelly
baselines in the main paper consequently employ its \texttt{torch} backend. A
natural concern is that a fused or Triton backend would be substantially faster,
thereby diminishing the apparent advantage of VD-TTFS. This appendix addresses the
concern directly. A fused multi-step LIF network---executed with SpikingJelly's
accelerated \texttt{cupy} and \texttt{triton} backends---nonetheless evaluates
\emph{every} one of the $T$ timesteps, performing neither early exit nor
event-sparsity exploitation. Its latency therefore constitutes a lower bound on the
cost of any fused TTFS backend (including a hypothetical TTFS-Triton kernel) on the
same network: should VD-TTFS prove faster than this bound, the improvement is
attributable to the algorithm (temporal collapse, arrival-ordered early exit, and
event sparsity) rather than to a deficiency of the framework.

\subsection{Experimental setup}\label{app:sjfused-setup}
We instantiate spiking networks with the identical layer structure and identical
number of timesteps $T$ as Task~1, Task~2, and Task~3, substituting standard LIF
neurons and random weights and inputs; inference accuracy is immaterial, as only
latency and memory are of interest. Each network is executed in multi-step mode
under the \texttt{triton} neuron backend, which is available in SpikingJelly master
(version 0.0.0.0.15 and later); the released 0.0.0.0.14 exposes only the
\texttt{torch} and \texttt{cupy} backends. All measurements are performed on a
single NVIDIA RTX~5070~Ti (16\,GB) with SpikingJelly~0.0.0.0.15, PyTorch~2.12 /
CUDA~13, and Triton~3.7. Peak GPU memory is measured with \texttt{nvidia-smi}
(per-process, allocator-agnostic, inclusive of the CUDA context);
\texttt{torch.cuda.max\_memory\_allocated} is not used, as it omits allocations
performed outside the PyTorch caching allocator, including those of the
\texttt{cupy} and \texttt{triton} backends.

\subsection{Acceptance threshold}\label{app:sjfused-threshold}
We adopt as the reference the per-sample inference latency of VD-TTFS on the same
device, obtained from the cross-platform measurements: $4.9\,\mu$s (Task~1),
$274\,\mu$s (Task~2), and $899\,\mu$s (Task~3) on the RTX~5070~Ti. A fused
multi-step configuration is deemed non-competitive if its per-sample latency
exceeds ten times the corresponding VD-TTFS latency, or if it cannot be executed
within the memory of the device. A configuration is classified \emph{infeasible}
when it does not terminate within the resulting latency budget (a factor of ten,
i.e.\ $0.49$\,s, $27.4$\,s, and $9.69$\,s for the respective full workloads) or
raises an out-of-memory condition.

\subsection{Results}\label{app:sjfused-results}
Table~\ref{tab:sjfused} reports, per task, the VD-TTFS reference latency, the
Triton multi-step latency, and the resulting determination. A dash denotes that
no latency could be obtained because the configuration is infeasible under the
criterion of Section~\ref{app:sjfused-threshold}.

\begin{table}[t]
\centering
\caption{SpikingJelly multi-step LIF (\texttt{triton} backend), with structure and
$T$ identical to each task and random weights, against the VD-TTFS reference on the
same RTX~5070~Ti. Latencies are per sample.}
\label{tab:sjfused}
\begin{tabular}{lrrl}
\toprule
Task ($T$) & VD-TTFS (ms/sample) & Triton (ms/sample) & Determination \\
\midrule
Task 1 ($T{=}80$)  & 0.0049 & 0.049 (peak 13{,}468\,MiB) & $10.0\times$ slower \\
Task 2 ($T{=}680$) & 0.274  & ---   & infeasible (memory) \\
Task 3 ($T{=}160$) & 0.899  & ---   & infeasible (latency) \\
\bottomrule
\end{tabular}
\end{table}

\subsection{Findings}\label{app:sjfused-find}
\begin{enumerate}
  \item \textbf{Backend invariance.} On Task~1 the \texttt{triton}, \texttt{cupy},
    and \texttt{torch} backends yield an identical per-sample latency
    ($0.049$\,ms) and, when measured with \texttt{nvidia-smi}, an identical peak
    memory ($\approx13.5$\,GB). As the elementwise neuronal update constitutes a
    negligible fraction of the per-timestep cost, the choice of backend influences
    neither throughput nor memory.
  \item \textbf{Latency lower bound (Task~1).} The Triton multi-step latency
    ($0.049$\,ms/sample) is exactly $10.0\times$ that of VD-TTFS
    ($0.0049$\,ms/sample). The disparity is structural rather than
    implementational: the multi-step pass evaluates all $80$ timesteps, whereas
    VD-TTFS collapses the temporal dimension and terminates at the first spike. A
    hypothetical TTFS-Triton kernel, subject to the same exhaustive traversal,
    would incur the same bound.
  \item \textbf{Memory infeasibility (Task~2).} For VGG-16 at $T{=}680$, multi-step
    execution simultaneously materializes the activations of all $680$ timesteps.
    The resulting footprint exhausts the $16$\,GB capacity of the device, raising an
    out-of-memory condition during execution; under memory configurations that
    avert the immediate allocation failure, the forward pass does not terminate
    within the latency budget. The configuration is therefore infeasible on this
    hardware.
  \item \textbf{Latency infeasibility (Task~3).} For the DVSGesture network at
    $T{=}160$, the Triton multi-step forward does not terminate within ten times
    the VD-TTFS latency; the $9.69$\,s budget is exceeded by more than a factor of
    forty without completion.
  \item \textbf{Conclusion.} The inability of the most heavily optimized
    SpikingJelly backend to match VD-TTFS---by latency on Task~1 and by feasibility
    on Task~2 and Task~3---confirms that the efficiency of VD-TTFS originates in the
    algorithm rather than in framework or backend implementation. That a fused
    Triton execution completes only one of the three workloads, and is an order of
    magnitude slower on that one, is itself evidence of this advantage.
\end{enumerate}

\subsection{Soundness of the single-step baseline}\label{app:sjfused-singlestep}
The SpikingJelly baselines reported in the main text employ the \emph{single-step}
execution mode, and complete every workload they are applied to: on the
RTX~5070~Ti, Task~1 in $0.996$\,s (peak $750$\,MiB) and Task~2 in $29.99$\,s (peak
$3876$\,MiB); on the datacenter V100 of the main text, $1.81$\,s and $38.4$\,s,
respectively. This stands in deliberate contrast to the multi-step configuration of
this appendix, which is infeasible for Task~2 and Task~3
(Section~\ref{app:sjfused-find}).

The contrast is explained by memory complexity. Single-step execution advances the
network one timestep at a time, retaining only the activations of the current
timestep together with the persistent neuronal state (the membrane potential); its
activation memory is therefore $\mathcal{O}(N)$ and \emph{independent of $T$}.
Multi-step execution instead materializes the activations of all $T$ timesteps at
once, giving $\mathcal{O}(N\,T)$ memory. For VGG-16 at $T{=}680$ the single-step
pass occupies $3.9$\,GB and fits comfortably, whereas the multi-step pass exceeds
the $16$\,GB device. Single-step execution further avoids the fused-kernel
autotuning over large per-timestep tensors that fails to terminate in the
multi-step case. The two modes thus embody a memory--throughput trade-off:
single-step incurs per-timestep launch overhead but is memory-bounded and
reliably terminates, whereas multi-step amortizes that overhead at the cost of
$\mathcal{O}(N\,T)$ memory that is prohibitive at long $T$.

These observations establish that the baseline choice is sound. Single-step is not
a deliberately weakened configuration: it is the only SpikingJelly mode that
executes the full TTFS dynamics of the deep, long-window networks within commodity
memory. The ostensibly faster fused/Triton multi-step alternative is, on identical
hardware, either no faster (Task~1, where the two modes share the same latency) or
infeasible (Task~2 and Task~3). The single-step baseline therefore represents the
best achievable SpikingJelly performance on these workloads and constitutes a fair,
indeed favorable, point of comparison. Even so, it remains an order of magnitude
slower than VD-TTFS---$20.3\times$ on Task~1 and $10.9\times$ on Task~2 on the
RTX~5070~Ti---a margin that, by the preceding analysis, a fused backend cannot
recover.

\section{Justification of the Evaluation Suite}\label{app:tasks}
The main text evaluates three tasks---LeNet on MNIST ($T{=}80$), VGG-16 on
CIFAR-10 ($T{=}680$), and a seven-layer recurrent network on the event-based
DVSGesture stream ($T{=}160$)---spanning shallow to deep architectures, static to
event-based data, and short to long temporal windows. We argue that the scale of
this suite is appropriate.

Increasing the network depth or the input complexity raises both the per-timestep
cost and, decisively, the memory footprint of any timestep-parallel (multi-step)
execution, whose activation memory scales as $\mathcal{O}(N\,T)$
(Section~\ref{app:sjfused-singlestep}). The de facto high-performance execution
path for spiking networks---the fused multi-step (CuPy/Triton) backend of
SpikingJelly, the most widely adopted framework---is already at its feasibility
limit on the present suite: as established in Section~\ref{app:sjfused}, it
exhausts the $16$\,GB device on VGG-16/CIFAR-10 and fails to terminate on the
DVSGesture network. A strictly larger network or a higher-resolution dataset would
place this baseline even further beyond reach; it could not be executed at all, and
hence could not serve as a point of comparison.

The remaining, feasible means of executing the full TTFS dynamics---SpikingJelly in
single-step mode and a hand-written step-by-step CUDA implementation---both
traverse every one of the $T$ timesteps without early exit or sparsity
exploitation, and are already markedly slower than VD-TTFS at the present scale:
the single-step baseline by $10.9\times$--$20.3\times$
(Section~\ref{app:sjfused-singlestep}) and the step-by-step CUDA baseline by up to
$36.8\times$ (main text). As both costs grow with network size and window length,
this margin would only widen on larger tasks.

The chosen suite therefore occupies the appropriate operating point: it is large
enough that the standard accelerated SNN baseline has reached the limit of
feasibility, yet remains tractable for the slower feasible baselines against which
VD-TTFS is measured. Larger benchmarks would not alter the qualitative
conclusion---they would merely preclude the baselines themselves---so the
three-task design of the main text is well justified.

\end{document}
